\documentclass[10pt,a4paper]{amsart}
\usepackage[margin=19mm]{geometry}
\usepackage{amsmath,amssymb,mathtools}
\usepackage[scr=rsfs,cal=euler]{mathalfa}
\usepackage{xfrac}
\usepackage[unicode,hidelinks,bookmarks=false]{hyperref}
\newcommand{\ie}{i.e.\ }
\newcommand{\cf}{cf.\ }
\newcommand{\resp}{respectively\ }
\newcommand{\Subsection}{\subsection}
\providecommand{\nobreakdash}{\nobreak\hbox{-}}
\setlength{\emergencystretch}{2em}
\multlinegap0pt
\usepackage{amssymb}
\newtheorem{thm}{Theorem}
\newtheorem{lem}{Lemma}
\title{The Gibbs phenomenon for the Krawtchouk polynomials}
\author{John Cullinan, Elisabeth Young}
\date{Source: arXiv:2603.05408v1 (CC BY 4.0)}
\begin{document}
\maketitle

\noindent\emph{Source and licence.} The Gibbs phenomenon for the Krawtchouk polynomials. Version arXiv:2603.05408v1 (CC BY 4.0), distributed by arXiv under the Creative Commons Attribution 4.0 International licence, http://creativecommons.org/licenses/by/4.0/, as recorded in the arXiv licence metadata for this exact version and shown on the abstract page https://arxiv.org/abs/2603.05408v1. Full licence notice: \texttt{licenses/arxiv-2603.05408-CC-BY-4.0.txt}.\par\smallskip

\noindent\textit{Editorial scope.} Counted unit: the article's thm induction\_theorem (source0.tex lines 425--427). The article's complete original proof of it is delivered with it (source0.tex lines 504--575, one \texttt{proof} environment of 72 lines, which the article places away from the statement). No mathematics is written, repaired or re-derived: the statement and the proof are the article's own lines, in the article's own notation, and no retained line is substituted or re-flowed.  Retained uncounted as the source-local context the proof needs: lines 416--418; lines 420--422; lines 430--433; lines 435--438; lines 444--446; lines 449--454.  Nothing was omitted from any retained span, so the package carries no omission record.  No retained line contains a citation, so the wrapper carries no bibliography and no bibliography entry.  No retained line contains a figure environment, an image-inclusion command or drawing code, so no image is delivered and the package declares no asset dependency.  Editorial formatting only, in addition: an AMS \texttt{amsart} preamble that restates the article's own declarations and macros used by the retained lines, namely the environments \texttt{thm}, \texttt{lem} on separate counters, together with the standard packages those lines need.  The retained environments print in this document as Theorem 1, Lemma 1 in order of appearance, whereas the article prints the same items with the numbering of its own full document.  The complete native source of the article as distributed in the arXiv e-print for this exact version is bundled unchanged as \texttt{sources/195801/source0.tex}.
\begin{align} \label{S(M)_definition}
S(M) = \frac{1}{2^{2M-1}}\binom{2M}{M} \sum_{m=0}^{M-1} (-1)^m \frac{\binom{M}{m}}{\binom{2M}{2m}} \sum_{\ell = 0}^{m} \frac{(-1)^{\ell}\binom{M}{\ell}}{2m+1-2\ell}.
\end{align}

\begin{align} \label{T(M)_definition}
T(M) = 2 \sum_{k=1}^{M} \frac{1}{M+k},
\end{align}

\begin{align}
C(m,M) &=(-1)^m \frac{\binom{M}{m}}{\binom{2M}{2m}} \\
D(m,M) &= C(m,M) \left(\frac{2M-2m+1}{2M+2} \right). \label{D(m,M)_definition}
\end{align}

\begin{align}
\left( \frac{2M+1}{2M+2} \right) C(m,M+1) &= D(m,M),\text{ and} \label{CD_identity_1} \\
D(m,M) - D(m+1,M) &= C(m,M). \label{CD_identity_2} 
\end{align}

\begin{align} \label{X(m,M)_identity}
X(m,M+1) = X(m,M) - X(m-1,M).
\end{align}

\begin{lem} \label{X(M,M+1)_lemma}
With all notation as above, we have 
\[
X(M,M+1) = (-1)^{M+1}\left(1 - \frac{2^{2M+1}}{\binom{2M+1}{M+1}}\right).
\]
\end{lem}

\begin{thm} \label{induction_theorem}
For $M \geq 1$ define $S(M)$ and $T(M)$ as in Equations (\ref{S(M)_definition}) and (\ref{T(M)_definition}).  Then $S(M) = T(M)$ for all $M \geq 1$.
\end{thm}

\begin{proof}[Proof of Theorem \ref{induction_theorem}]
We verify $S(1) = T(1) =1$.  Fix $M \geq 1$, assume $S(M) = T(M)$, and consider $S(M+1)$:
\begin{align*}
S(M+1) &= \frac{1}{2^{2M+1}} \binom{2M+2}{M+1}\sum_{m=0}^M C(m,M+1)X(m,M+1) \\
&=  \frac{1}{2^{2M-1}} \binom{2M}{M}  \left(\frac{2M+1}{2M+2}\right) \sum_{m=0}^M C(m,M+1)X(m,M+1).
\end{align*}
Applying Equation (\ref{CD_identity_1}) we  rewrite $S(M+1)$ as
\begin{align}\label{S(M+1)_rewritten}
S(M+1) = \frac{1}{2^{2M-1}} \binom{2M}{M} \sum_{m=0}^{M} D(m,M)X(m,M+1).
\end{align}
Next, write
\[
\sum_{m=0}^{M} D(m,M)X(m,M+1) =\sum_{m=0}^{M-1} D(m,M)X(m,M+1) + D(M,M)X(M,M+1).
\]
By direct substitution in Equation (\ref{D(m,M)_definition}), we have 
\begin{align} \label{D(M,M)}
D(M,M) = \frac{(-1)^M}{2M+2}.
\end{align}
Applying Lemma \ref{X(M,M+1)_lemma}, we then have 
\begin{align} \label{D(M,M)X(M,M+1)}
\frac{\binom{2M}{M}}{2^{2M-1}}D(M,M)X(M,M+1) = \frac{2}{2M+1} - \frac{\binom{2M}{M}}{2^{2M-1}(2M+2)}.
\end{align}
Now use Equation (\ref{X(m,M)_identity}) to write
\begin{align*}
\sum_{m=0}^{M-1} D(m,M)X(m,M+1) &= \sum_{m=0}^{M-1} D(m,M)\left(X(m,M) - X(m-1,M)\right) \\
&=\sum_{m=0}^{M-1} D(m,M)X(m,M) - \sum_{m=0}^{M-1} D(m,M)X(m-1,M).
\end{align*}
Since $X(-1,M) = 0$ for all $M$, we reindex the second sum as
\begin{align} \label{reindex}
\sum_{m=0}^{M-1} D(m,M)X(m-1,M) = \sum_{m=0}^{M-2} D(m+1,M)X(m,M).
\end{align}
Now we write
\begin{align} \label{breakup}
\sum_{m=0}^{M-1} D(m,M)X(m,M) = \sum_{m=0}^{M-2} D(m,M)X(m,M) + D(M-1,M)X(M-1,M).
\end{align}
Applying Equations (\ref{reindex}) and (\ref{breakup}) we have
\begin{align*}
\sum_{m=0}^{M-1} D(m,M)X(m,M+1) &= \sum_{m=0}^{M-2}\left(D(m,M) - D(m+1,M)\right)X(m,M) \\
& + D(M-1,M)X(M-1,M).
\end{align*}
By Equation (\ref{CD_identity_2}), we have 
\begin{align} \label{condense}
\sum_{m=0}^{M-2}\left(D(m,M) - D(m+1,M)\right)X(m,M) = \sum_{m=0}^{M-2}C(m,M)X(m,M),
\end{align}
and by direct calculation (using Lemma \ref{X(M,M+1)_lemma}) we verify that 
\begin{align} \label{simplification}
\frac{\binom{2M}{M}}{2^{2M-1}}D(M-1,M)X(M-1,M) &= \frac{2}{2M-1} - \frac{2}{2M+2} \\
&- \frac{3\binom{2M}{M}}{2^{2M-1}(2M-1)(2M+2)}. \nonumber
\end{align}
Multiplying Equation (\ref{condense}) by $\frac{\binom{2M}{M}}{2^{2M-1}}$ and adding the right-hand-sides of Equations (\ref{D(M,M)X(M,M+1)}) and  (\ref{simplification}), we have
\begin{align*}
S(M+1) &= \frac{\binom{2M}{M}}{2^{2M-1}} \sum_{m=0}^{M-2}C(m,M)X(m,M)  - \frac{\binom{2M}{M}}{2^{2M-1}(2M+2)}\left(1 + \frac{3}{2M-1} \right) \\&+\frac{2}{2M-1} - \frac{2}{2M+2} + \frac{2}{2M+1}.
\end{align*}
Since 
\begin{align*}
\frac{\binom{2M}{M}}{2^{2M-1}(2M+2)}\left(1 + \frac{3}{2M-1} \right) &= \frac{\binom{2M}{M}}{2^{2M-1}(2M-1)}\\
\end{align*}
and
\begin{align*}
-\frac{\binom{2M}{M}}{2^{2M-1}(2M-1)} + \frac{2}{2M-1} = \frac{1}{2M-1} \left(\frac{2^{2M}}{\binom{2M}{M}} -1 \right) = C(M-1,M)X(M-1,M) 
\end{align*}
(using Lemma \ref{X(M,M+1)_lemma}, again), we have
\begin{align*}
S(M+1) &= S(M)  - \frac{1}{M+1} + \frac{2}{2M+1}.
\end{align*}
Applying the induction hypothesis $S(M) = T(M)$, we have
\begin{align*}
S(M+1) &= 2 \sum_{k=1}^M \frac{1}{M+k} - \frac{1}{M+1} + \frac{2}{2M+1} \\
&=2\sum_{k=1}^{M+1} \frac{1}{M+1+k} = T(M+1),
\end{align*}
which completes the proof.
\end{proof}

\end{document}
